%! Author  = Omar Iskandarani
%! Date    = June 2026
%! ORCID   = 0009-0006-1686-3961
%! Target  = Physical Review A / European Physical Journal Plus

% ============================================================
%  Closure Requirements for a Compton-Core Vortex-Filament Scaling Model:
%  What Is Fixed Algebraically and What Requires New Dynamics
% ============================================================

\documentclass[aps,pra,preprint,longbibliography,nofootinbib]{revtex4-2}

\usepackage{amsmath,amssymb,amsthm,mathtools}
\usepackage{booktabs}
\usepackage{xcolor}
\usepackage{enumitem}
\usepackage{siunitx}
\usepackage[colorlinks=true,
            linkcolor=blue!60!black,
            citecolor=blue!60!black,
            urlcolor=blue!60!black]{hyperref}
\usepackage{cleveref}
\usepackage{microtype}

% ── Theorem environments ─────────────────────────────────────
\newtheorem{postulate}{Postulate}
\newtheorem{lemma}{Lemma}
\newtheorem{proposition}{Proposition}
\newtheorem{corollary}{Corollary}
\newtheorem{remark}{Remark}
\newtheorem{closure}{Closure}

% ── Notation ─────────────────────────────────────────────────
% Circulatory arrow subscript (reads as "circulatory flow")
\newcommand{\rotarrow}{%
  \mathchoice{\mkern-2mu\scriptstyle\boldsymbol{\circlearrowleft}}%
             {\mkern-2mu\scriptstyle\boldsymbol{\circlearrowleft}}%
             {\mkern-2mu\scriptscriptstyle\boldsymbol{\circlearrowleft}}%
             {\mkern-2mu\scriptscriptstyle\boldsymbol{\circlearrowleft}}}

\newcommand{\vc}{\mathbf{v}_{\rotarrow}}        % circulatory velocity (vector)
\newcommand{\vnorm}{\lVert\mathbf{v}_{\!\boldsymbol{\circlearrowleft}}\rVert} % magnitude

\newcommand{\rhof}{\rho_{\!f}}                   % background fluid density
\newcommand{\rhoE}{\rho_{\!E}}                   % kinetic energy density
\newcommand{\rhom}{\rho_{\!m}}                   % mass-equivalent density
\newcommand{\rhocore}{\rho_{\mathrm{core}}}      % core density
\newcommand{\rc}{r_c}                            % core radius
\newcommand{\Stv}{S_t}                           % relativistic time-dilation factor
\newcommand{\omC}{\omega_C}                      % Compton angular frequency
\newcommand{\lc}{\lambda_c}                      % Compton wavelength
\newcommand{\Rinf}{R_\infty}                     % Rydberg constant
\newcommand{\Tc}{\mathcal{T}_c}                  % vortex line tension at core
\newcommand{\Fgr}{F_{\mathrm{gr}}}              % gravitational tension scale c^4/4G
\newcommand{\Gz}{\Gamma_0}                       % circulation quantum
% ─────────────────────────────────────────────────────────────

\begin{document}

\title{Closure Requirements for a Compton-Core Vortex-Filament Scaling Model:\\
       What Is Fixed Algebraically and What Requires New Dynamics}

\author{Omar Iskandarani}
\thanks{ORCID: \href{https://orcid.org/0009-0006-1686-3961}{0009-0006-1686-3961}}
\affiliation{Independent Researcher, Groningen, The Netherlands}
\email{info@omariskandarani.com}
\date{June 2026}

% ─────────────────────────────────────────────────────────────

\begin{abstract}
We analyze which parts of a Compton-core vortex-filament model of the
electron are fixed algebraically and which parts require additional
physical closure.  The core postulate identifies the angular rotation
rate of the core with the electron Compton frequency,
$\Omega_{\mathrm{core}}=\omega_C=m_e c^2/\hbar$.  This postulate alone
only gives $\lVert \mathbf{v}_{\!\boldsymbol{\circlearrowleft}}\rVert=\omega_C r_c$;
it does not determine the characteristic circulation speed, the
fine-structure constant, the line tension, or the core density.  To close
the model one must add: (i) a spectroscopic speed closure equivalent to
using the Rydberg constant or $\alpha$ as input; (ii) a core-diameter
energy normalization for the line tension; and (iii) a thin-wall
Laplace-pressure closure for the density.  With these assumptions the
scales
$\lVert \mathbf{v}_{\!\boldsymbol{\circlearrowleft}}\rVert=(\pi\hbar cR_\infty/m_e)^{1/2}$,
$r_c=\lVert \mathbf{v}_{\!\boldsymbol{\circlearrowleft}}\rVert/\omega_C=r_e/2$,
$\Gamma_0=2\pi r_c\lVert \mathbf{v}_{\!\boldsymbol{\circlearrowleft}}\rVert$,
$\mathcal{T}_c=m_ec^2/(2r_c)$, and
$\rho_{\mathrm{core}}=m_ec^2/(2\pi\lVert \mathbf{v}_{\!\boldsymbol{\circlearrowleft}}\rVert^2r_c^3)$
follow consistently.  We therefore do not claim a first-principles
derivation of $\alpha$ or of $\rho_{\mathrm{core}}$.  Instead we identify
the minimal missing ingredients required for such a derivation and give
phenomenological test templates in terms of a short-range potential
correction or an electron form factor.
\end{abstract}

\maketitle

% =============================================================
\section{Introduction}
\label{sec:intro}
% =============================================================

Localized vortex-filament models of particle-like excitations trace
their lineage to Helmholtz and Kelvin~\cite{Helmholtz1867,Kelvin1867}
and continue to appear in topological soliton theory~\cite{FaddeevNiemi1997},
analogue-gravity models~\cite{Barcelo2005}, and superfluid-vacuum
analogies~\cite{Volovik2003}.  Such models require at least a core
radius, a characteristic circulation speed, a circulation, a line
tension, and an effective core density.  The central issue is not
whether these quantities can be written in terms of known constants;
they can.  The issue is which relations are independent physical
postulates and which are merely algebraic rewritings of standard
spectroscopic data.

The present paper therefore adopts a deliberately conservative aim.  We
ask what follows from a Compton-core hypothesis, what additionally must
be supplied to close the model, and what would be required to turn the
closure into a first-principles derivation.  The result is a hierarchy of
assumptions rather than a claimed derivation of the fine-structure
constant.  In particular, the relation between the circulation speed and
$\alpha$ is treated as a spectroscopic speed closure equivalent to using
$R_\infty$ as input.  Likewise, the line tension and the density require
separate mechanical closure assumptions.

This distinction is essential.  The Compton-core postulate gives only
$\lVert \mathbf{v}_{\!\boldsymbol{\circlearrowleft}}\rVert=\omega_C r_c$.
It does not by itself fix either
$\lVert \mathbf{v}_{\!\boldsymbol{\circlearrowleft}}\rVert$ or $r_c$.
A genuine derivation of the speed would require an independent dynamical
calculation of the electromagnetic coupling, or equivalently of
$\alpha$, without importing $R_\infty$.  A genuine derivation of the line
tension would require a specified core-energy functional.  A genuine
derivation of the density would require a specified interface model or
stress tensor.  These requirements are made explicit below.

Throughout we use SI units, with $\hbar=h/(2\pi)$, $m_e$, $c$, and $e$
carrying their CODATA 2022 recommended meanings~\cite{MohrTaylor2022}.
CODATA 2022 is based on a least-squares adjustment of data available
through 31 December 2022~\cite{MohrTaylor2022}.  We denote the
circulatory velocity field by $\mathbf{v}_{\!\boldsymbol{\circlearrowleft}}$
and its magnitude by
$\vnorm=\lVert \mathbf{v}_{\!\boldsymbol{\circlearrowleft}}\rVert$.

% =============================================================
\section{The Compton-Core Postulate}
\label{sec:postulate}
% =============================================================

\begin{postulate}[Compton-core hypothesis]
\label{post:compton}
The vortex core undergoes solid-body rotation at the electron
Compton angular frequency,
\begin{equation}
  \Omega_{\mathrm{core}} = \omC \equiv \frac{m_e c^2}{\hbar}.
  \label{eq:postulate}
\end{equation}
The magnitude of the circulatory velocity at the core boundary
of radius $\rc$ is therefore
\begin{equation}
  \vnorm = \omC\,\rc.
  \label{eq:vc_rc_link}
\end{equation}
\end{postulate}

\begin{remark}[Physical motivation]
\label{rem:motivation}
For a solid-body rotating core of radius $\rc$ and angular
velocity $\Omega$, the tangential speed at the boundary is
$\vnorm = \Omega\,\rc$.
The Compton frequency $\omC = m_e c^2/\hbar$ is the natural
frequency associated with the electron rest energy: it is the
frequency at which one quantum $\hbar\omC = m_e c^2$ is stored.
Postulate~\ref{post:compton} is the simplest resonance condition
consistent with the electron rest energy as the core-energy scale,
and it provides a geometric link between the rotation rate and the
core boundary.
It is a hypothesis about the core dynamics, not a consequence of
standard quantum electrodynamics, and it is testable: any
independent measurement of $\rc$ inconsistent with
$\rc = \vnorm/\omC$ (where $\vnorm$ is determined by
Sec.~\ref{sec:derived}) falsifies the postulate.
\end{remark}

% =============================================================
\section{The Spectroscopic Speed Closure}
\label{sec:rydberg}
% =============================================================

\subsection{Equivalence of $R_\infty$ and $\alpha$ as inputs}

The Rydberg constant satisfies
\begin{equation}
  R_\infty = \frac{\alpha^2 m_e c}{2h}
           = \frac{\alpha^2 m_e c}{4\pi\hbar}.
  \label{eq:Rydberg_standard}
\end{equation}
Thus $R_\infty$ and $\alpha$ carry the same information once
$\{m_e,c,h\}$ are specified.  Using $R_\infty$ instead of $\alpha$ is
therefore not an independent derivation of $\alpha$; it is a convenient
spectroscopic parametrization of the same dimensionless input.

\begin{closure}[Spectroscopic speed closure]
\label{cl:speed}
The characteristic circulation speed is related to the electromagnetic
coupling by
\begin{equation}
  \tilde\alpha \equiv \frac{2\vnorm}{c},
  \qquad
  R_\infty = \frac{\tilde\alpha^2 m_e c}{2h}.
  \label{eq:speed_closure}
\end{equation}
Equivalently,
\begin{equation}
  R_\infty = \frac{\vnorm^2 m_e}{\pi\hbar c}
           = \frac{\vnorm^2\omega_C}{\pi c^3}.
  \label{eq:rydberg_core}
\end{equation}
\end{closure}

\begin{remark}[What would be needed to derive the speed]
\label{rem:derive_speed}
Closure~\ref{cl:speed} is the strongest non-geometric input in the
model.  A first-principles derivation would require a dynamical
calculation of the dimensionless coupling
\begin{equation}
  \tilde\alpha = \mathcal{F}[\mathcal{I}_1,\mathcal{I}_2,\ldots],
  \label{eq:alpha_from_dynamics_placeholder}
\end{equation}
where the invariants $\mathcal{I}_i$ are defined by an explicit
field-theoretic or hydrodynamic action.  Unless Eq.~\eqref{eq:alpha_from_dynamics_placeholder}
can be evaluated without using $R_\infty$, atomic spectra, or the
CODATA value of $\alpha$, the speed remains calibrated rather than
derived.
\end{remark}

% =============================================================
\section{Closure of the Characteristic Speed}
\label{sec:derivation}
% =============================================================

Using Closure~\ref{cl:speed},
\begin{equation}
  \Rinf = \frac{\vnorm^2\,\omC}{\pi c^3}
  \quad\Longrightarrow\quad
  \vnorm^2 = \frac{\pi\,\hbar\,c\,\Rinf}{m_e}.
  \label{eq:vc_solve}
\end{equation}

\begin{proposition}[Characteristic circulation speed]
\label{prop:vc}
Given Closure~\ref{cl:speed} and the measured $\Rinf$,
\begin{equation}
  \boxed{
    \vnorm = \sqrt{\frac{\pi\,\hbar\,c\,\Rinf}{m_e}}
  }
  \label{eq:vc}
\end{equation}
is the unique positive closure value for the tangential
circulation speed.
The inputs are $\{m_e, c, \hbar, \Rinf\}$; since
$\Rinf \equiv \alpha^2 m_e c/(2h)$, this is equivalent to using
$\alpha$ as input.  The result is therefore calibrated by spectroscopy,
not predicted independently.
\end{proposition}

\subsubsection*{Dimensional check}
\[
  \Bigl[\sqrt{\hbar\,c\,\Rinf/m_e}\Bigr]
  = \sqrt{
      \frac{\mathrm{J{\cdot}s\cdot m{\cdot}s^{-1}\cdot m^{-1}}}
           {\mathrm{kg}}
    }
  = \sqrt{\mathrm{m^2\,s^{-2}}} = \mathrm{m\,s^{-1}}. \checkmark
\]

\subsubsection*{Numerical value}
$\vnorm = \SI{1.093\,845\,631\,5e6}{\meter\per\second}$
\quad(from CODATA 2022~\cite{MohrTaylor2022}).

% =============================================================
\section{Derived Mechanical Scales}
\label{sec:derived}
% =============================================================

The following scales follow algebraically from
Proposition~\ref{prop:vc} and Postulate~\ref{post:compton},
each with a direct classical-hydrodynamic interpretation.

% ------- 4.1 Core radius ------------------------------------
\subsection{Core radius and its relation to the classical electron radius}
\label{sec:rc}

\begin{equation}
  \rc = \frac{\vnorm}{\omC} = \frac{\vnorm\,\hbar}{m_e c^2}
      = \SI{1.408\,970\,16e-15}{\meter}.
  \label{eq:rc}
\end{equation}

\begin{remark}[Relation to the classical electron radius]
\label{rem:re}
The classical electron radius is
\begin{equation}
  r_e = \frac{e^2}{4\pi\varepsilon_0 m_e c^2}
      = \frac{\alpha\hbar}{m_e c}
      = \SI{2.817\,940\,3e-15}{\meter}.
  \label{eq:re}
\end{equation}
From \eqref{eq:rc} and $\alpha = 2\vnorm/c$ (established below),
$\rc = \alpha\hbar/(2m_e c) = r_e/2$.
We emphasize that $\rc$ is an \emph{effective model parameter},
not a claim about the physical extent of the electron.
The classical radius $r_e$ itself plays an analogous role in QED
— appearing in Thomson cross-sections, Compton scattering, and
vacuum polarization — without implying a physical "size."
Experimental upper bounds on the electron's spatial extent
(from $g{-}2$ and high-energy scattering measurements,
$r \lesssim \SI{1e-18}{\meter}$~\cite{PDG2022}) constrain the
\emph{point-like} quantum-mechanical extent of the charge
distribution, not the scale at which an effective vortex-core
model is constructed.
For a toroidal idealization in which the tube centre lies at distance
$R = \rc$ from the symmetry axis and the tube has minor radius
$a = \rc$ (horn-torus geometry), the outer boundary lies at
$R+a=2\rc=r_e$.  This gives a geometric interpretation of the
factor-of-two relation, but it is not evidence that the electron charge
has a resolved radius $r_e$.
\end{remark}

% ------- 4.2 Fine-structure constant -------------------------
\subsection{Fine-structure constant as input equivalence}
\label{sec:alpha}

\begin{proposition}[Algebraic equivalence of $\alpha$ and $\vnorm$]
\label{prop:alpha}
The ratio $\tilde\alpha \equiv 2\vnorm/c$
satisfies $\tilde\alpha = \alpha$ (the CODATA fine-structure
constant) because Closure~\ref{cl:speed} was imposed.
This is not a derivation of $\alpha$ from independent physics
but the algebraic consequence of $\Rinf \equiv \alpha^2 m_e c/(2h)$
combined with \eqref{eq:vc}.
\end{proposition}

\begin{proof}
$\tilde\alpha = (2/c)(\pi\hbar c\Rinf/m_e)^{1/2}$.
Substituting $\Rinf = \alpha^2 m_e c/(4\pi\hbar)$:
$\tilde\alpha = (2/c)(\alpha^2 c^2/4)^{1/2} = \alpha$. \qed
\end{proof}

Numerically, $\tilde\alpha = \num{7.297\,352\,567\,1e-3}$
versus CODATA $\alpha = \num{7.297\,352\,569\,3e-3}$;
the residual $3\times10^{-10}$ is the propagated uncertainty
from $\Rinf$ (relative uncertainty \SI{1.9e-12}{}) combined with
the rounding of $\{m_e, c, \hbar\}$.
It is a precision-consistency check, not a prediction.

% ------- 4.3 Circulation quantum ----------------------------
\subsection{Circulation quantum}
\label{sec:Gamma0}

The line integral of velocity around the core boundary defines the
circulation,
\begin{equation}
  \Gz \equiv \oint \vc\cdot d\mathbf{l}
           = 2\pi\,\rc\,\vnorm
           = \frac{2\pi\vnorm^2}{\omC}
           = \SI{9.683\,619e-9}{\meter\squared\per\second}.
  \label{eq:Gamma0}
\end{equation}
The ratio $\Gz/\kappa = \alpha^2/4 \approx \num{1.33e-5}$,
where $\kappa = h/m_e$ is the Onsager-Feynman quantum~\cite{Onsager1949,Feynman1955},
showing that the present circulation is smaller by $\alpha^2/4$.

% ------- 4.4 Vortex line tension ---------------------------
\subsection{Vortex line tension}
\label{sec:Tc}

The \emph{vortex line tension} is the energy stored in the core per
unit filament length~\cite{Saffman1992}.  It is not fixed by
Postulate~\ref{post:compton}.  We therefore introduce a separate
normalization closure.

\begin{closure}[Core-diameter energy normalization]
\label{cl:tension}
One electron rest-energy quantum is assigned to a filament segment of
length $2\rc$:
\begin{equation}
  \Tc \equiv \frac{m_e c^2}{2\rc} = \frac{\hbar\omC}{2\rc}.
  \label{eq:Tc_def}
\end{equation}
\end{closure}

The physical content of Closure~\ref{cl:tension} is that the core stores
one quantum of Compton energy $m_e c^2=\hbar\omega_C$ per core-diameter
length.  A first-principles derivation would require an explicit
energy functional whose on-shell line energy equals Eq.~\eqref{eq:Tc_def}.

\begin{corollary}[Line tension in fundamental constants]
\label{cor:Tc}
\begin{equation}
  \Tc = \frac{m_e c^2}{2\rc}
      = \frac{m_e^2 c^3}{\tilde\alpha\,\hbar}
      = \SI{29.053\,51}{\newton}.
  \label{eq:Tc}
\end{equation}
\end{corollary}

\begin{remark}[Disambiguation from the gravitational tension scale]
\label{rem:Fgr}
The quantity $\Fgr \equiv c^4/(4G)$ defines the maximum force
transmissible without forming a causal horizon in general
relativity~\cite{Gibbons2002,Schiller2006}.
It and $\Tc$ share the dimension of force (energy per unit length),
but are conceptually distinct:
$\Tc$ is a mechanical property of the vortex core, while $\Fgr$ is
a causal upper bound from spacetime physics.
Numerically $\Tc \approx \SI{29}{\newton}$ versus
$\Fgr \approx \SI{3e43}{\newton}$.
Their ratio
\begin{equation}
  \frac{\Tc}{\Fgr} = \frac{4\alpha_g}{\tilde\alpha}
                   \approx \num{9.60e-43},
  \label{eq:tension_ratio}
\end{equation}
where $\alpha_g = Gm_e^2/(\hbar c) \approx \num{1.752e-45}$,
involves only fundamental constants and has no free parameters.
\end{remark}

% ------- 4.5 Core mass density --------------------------------
\subsection{Core mass density from Laplace-pressure balance}
\label{sec:rhocore}

The density is also not fixed by the Compton-core postulate.  We close
it with a thin-wall pressure-balance model.

\begin{closure}[Thin-wall Laplace-pressure closure]
\label{cl:laplace}
For a cylindrical core of radius $\rc$ and line tension $\Tc$, define an
effective wall tension
\begin{equation}
  \gamma \equiv \frac{\Tc}{2\pi\rc}.
  \label{eq:wall_tension}
\end{equation}
The corresponding cylindrical Laplace pressure is
\begin{equation}
  P_L = \frac{\gamma}{\rc} = \frac{\Tc}{2\pi\rc^2}.
  \label{eq:Laplace}
\end{equation}
We then impose the mechanical balance
\begin{equation}
  \frac{1}{2}\rhocore\vnorm^2 = \frac{\Tc}{2\pi\rc^2}.
  \label{eq:pressure_balance}
\end{equation}
\end{closure}

Closure~\ref{cl:laplace} is a thin-interface assumption.  It is not a
generic theorem of vortex dynamics; other core stress tensors or outer
cutoffs would change the numerical coefficient.

\begin{corollary}[Core density]
\label{cor:rho}
Solving~\eqref{eq:pressure_balance} and substituting
\eqref{eq:Tc_def}:
\begin{equation}
  \boxed{
    \rhocore = \frac{\Tc}{\pi\rc^2\vnorm^2}
             = \frac{m_e c^2}{2\pi\,\vnorm^2\,\rc^3}.
  }
  \label{eq:rhocore}
\end{equation}
Numerically, $\rhocore = \SI{3.893\,4e18}{\kilogram\per\cubic\meter}$.
\end{corollary}

\subsubsection*{Dimensional check}
$[m_e c^2/(\vnorm^2\rc^3)]
= \mathrm{kg\,m^2\,s^{-2}\,/\,(m^2\,s^{-2}\cdot m^3)}
= \mathrm{kg\,m^{-3}}$. \checkmark

\begin{remark}[Status of the density result]
\label{rem:noncircular}
Equation~\eqref{eq:rhocore} is non-circular only because
Closure~\ref{cl:tension} fixes $\Tc$ before Closure~\ref{cl:laplace} is
applied.  It is nevertheless not a first-principles density derivation.
It is the density implied by the chosen tension normalization and the
chosen thin-wall pressure law.

For comparison, the usual local-induction estimate for the line energy
of a vortex filament contains a circulation, a density, and typically a
logarithmic dependence on an outer length scale~\cite{Saffman1992}.
Therefore, any future derivation of $\rhocore$ from a hydrodynamic action
must specify the core profile, the outer cutoff or renormalization rule,
and the stress tensor.  Without those ingredients, Eq.~\eqref{eq:rhocore}
should be read as a closure consequence, not as an independent prediction.
\end{remark}

% ------- 4.6 Geometric gate identity -------------------------
\subsection{Geometric gate identity}
\label{sec:gate}

\begin{lemma}[Geometric gate]
\label{lem:gate}
\begin{equation}
  \frac{\lc}{\pi\rc} = \frac{4}{\tilde\alpha} \approx 548.14,
  \label{eq:gate}
\end{equation}
where $\lc = h/(m_ec)$ is the electron Compton wavelength.
\end{lemma}

\begin{proof}
$\lc = 2\pi\hbar/(m_ec)$; $\rc = \tilde\alpha\hbar/(2m_ec)$:
$\lc/(\pi\rc) = 4/\tilde\alpha$. \qed
\end{proof}

This dimensionless ratio is pure algebra; it is a geometric property
of the length scales $\lc$ and $\rc$ and does not encode additional
physical content.

% =============================================================
\section{Derivation Chain and Hydrodynamic Interpretation}
\label{sec:chain}
% =============================================================

The complete derivation chain is:
\begin{equation}
  \underbrace{m_e,\;c,\;\hbar,\;R_\infty}_{\text{standard inputs}}
  \;\xrightarrow{\;\text{Closure\;\ref{cl:speed}}\;}
  \vnorm
  \;\xrightarrow{\;\text{Post.\;\ref{post:compton}}\;}
  \rc
  \xrightarrow{}
  \Gz
  \;\xrightarrow{\;\text{Closure\;\ref{cl:tension}}\;}
  \Tc
  \;\xrightarrow{\;\text{Closure\;\ref{cl:laplace}}\;}
  \rhocore.
  \label{eq:chain}
\end{equation}
The classical-hydrodynamic interpretation of each scale is:

\medskip
\begin{tabular}{@{}ll@{}}
$\vnorm = \omC\rc$ & tangential speed at solid-body core boundary \\
$\rc = \vnorm/\omC$ & core radius from solid-body rotation \\
$\Gz = 2\pi\rc\vnorm$ & circulation (line integral of velocity) \\
$\Tc = m_e c^2/(2\rc)$ & line tension from core-diameter energy closure \\
$\rhocore$ & core density from thin-wall pressure-balance closure \\
\end{tabular}

% =============================================================
\section{Status Table}
\label{sec:status}
% =============================================================

\begin{table}[h]
\centering
\caption{Epistemic status of all quantities.  ``Closure'' denotes an
explicit model assumption beyond the Compton-core postulate; ``Algebraic''
denotes a consequence of earlier rows.}
\label{tab:status}
\begin{tabular}{llll}
\toprule
Quantity & Symbol & Status & Note \\
\midrule
Compton frequency  & $\omC = m_e c^2/\hbar$     & Standard def.  & \cite{MohrTaylor2022} \\
Core rotation rate & $\Omega_{\rm core}=\omC$   & Postulate       & Post.~\ref{post:compton} \\
Rydberg/FSC input  & $\Rinf\equiv\alpha$        & Empirical input & \eqref{eq:Rydberg_standard} \\
Speed relation     & $2\vnorm/c=\tilde\alpha$   & Closure         & Closure~\ref{cl:speed} \\
Circ.\ speed       & $\vnorm$                   & Calibrated      & Prop.~\ref{prop:vc} \\
Core radius        & $\rc = r_e/2$              & Post.+closure   & \eqref{eq:rc} \\
Circulation        & $\Gz$                      & Algebraic       & \eqref{eq:Gamma0} \\
Line tension       & $\Tc$                      & Closure         & Closure~\ref{cl:tension} \\
Core density       & $\rhocore$                 & Closure result  & Closure~\ref{cl:laplace} \\
Geom.\ gate        & $\lc/(\pi\rc)$             & Algebraic       & Lemma~\ref{lem:gate} \\
Background density & $\rhof$                    & Open            & needs independent coupling \\
Topological state  & knot type $K$              & Open            & not determined here \\
\bottomrule
\end{tabular}
\end{table}

% =============================================================
\section{What Would Be Required for a First-Principles Derivation}
\label{sec:requirements}
% =============================================================

The preceding sections define a consistent closure hierarchy.  A full
derivation would require replacing each closure by a dynamical equation.
The minimal requirements are as follows.

\subsection{Independent derivation of the speed or fine-structure constant}

One must specify an action or Hamiltonian containing both the
circulatory degrees of freedom and the electromagnetic coupling, for
example schematically
\begin{equation}
  S = \int d^4x\,\mathcal{L}\bigl[\mathbf{v}_{\!\boldsymbol{\circlearrowleft}},
  \rho, A_\mu, \psi;\Pi_i\bigr],
  \label{eq:needed_action}
\end{equation}
with dimensionless parameters or topological invariants $\Pi_i$.  The
model must then predict
\begin{equation}
  \alpha = \mathcal{F}(\Pi_i),
  \qquad
  \vnorm = \frac{\alpha c}{2},
  \label{eq:needed_alpha}
\end{equation}
without using $R_\infty$, atomic spectra, or the CODATA value of
$\alpha$.  This is the missing step if the goal is to derive the speed
rather than calibrate it.

\subsection{Independent derivation of the line tension}

A line tension must follow from an energy functional.  For a cylindrical
core profile this would require an expression of the form
\begin{equation}
  \mathcal{T}_{\rm dyn}
  = \int_0^{R_{\rm out}}2\pi r\,dr\,
    \left[\frac{1}{2}\rho(r)v_\theta(r)^2
    +U(\rho,\nabla\rho)+\mathcal{E}_{\rm int}(r)\right],
  \label{eq:needed_tension}
\end{equation}
plus a renormalization prescription for the outer scale $R_{\rm out}$.
Closure~\ref{cl:tension} is derived only if
\begin{equation}
  \mathcal{T}_{\rm dyn}=\frac{m_ec^2}{2r_c}
  \label{eq:tension_target}
\end{equation}
follows as the stationary on-shell value.

\subsection{Independent derivation of the density}

The density requires a stress-balance calculation.  A sufficient route is
to define a radial stress tensor $\sigma_{ij}$ and show that the normal
stress jump at the core wall gives
\begin{equation}
  \Delta P = \sigma_{rr}^{\rm(out)}-\sigma_{rr}^{\rm(in)}
           = \frac{\gamma}{r_c},
  \qquad
  \gamma = \frac{\partial E}{\partial A},
  \label{eq:needed_laplace}
\end{equation}
with $\gamma=\mathcal{T}_c/(2\pi r_c)$ as a consequence rather than a
definition.  Only then does $\rho_{\mathrm{core}}$ become dynamically
predicted.

\subsection{Independent low-energy observable}

Finally, the model must produce an observable correction.  A generic
low-energy parametrization is
\begin{equation}
  F_e(q^2)=1-\frac{1}{6}\langle r^2\rangle_{\rm eff}q^2+\mathcal{O}(q^4),
  \qquad
  \langle r^2\rangle_{\rm eff}=\chi r_c^2,
  \label{eq:form_factor_template}
\end{equation}
where $\chi$ must be predicted.  Equivalently, the theory may predict a
short-range potential correction
\begin{equation}
  \Delta V_\eta(r)=
  \eta\,\frac{e^2}{4\pi\varepsilon_0 r}\,f(r/r_c),
  \qquad f(0)\;\mathrm{finite},\quad f(x\gg1)\to0,
  \label{eq:potential_template}
\end{equation}
where $\eta$ and $f$ must be fixed by the model.  Without such a
coefficient or shape function, the core scale is not experimentally
sharp.

% =============================================================
\section{Phenomenological Tests and Falsifiers}
\label{sec:falsifiers}
% =============================================================

The present closure model does not yet make a parameter-free prediction
for atomic spectra.  It does, however, define test templates that become
falsifiable once the coefficients in Eqs.~\eqref{eq:form_factor_template}
and~\eqref{eq:potential_template} are fixed.

\medskip
\noindent\textbf{(i) Spectroscopic potential test.}
For any specified $\Delta V_\eta(r)$, hydrogenic energy shifts are
\begin{equation}
  \Delta E_{n\ell}
  = \int_0^\infty |R_{n\ell}(r)|^2\,\Delta V_\eta(r)\,r^2\,dr.
  \label{eq:lamb_template}
\end{equation}
Existing bound-state QED calculations already account for the Lamb
shift to high precision~\cite{Bethe1947,Karshenboim2005}.  Therefore,
any order-unity $\eta$ producing an unobserved shift is excluded.  If the
model predicts $\eta\ll1$, the corresponding bound must be stated.

\medskip
\noindent\textbf{(ii) Form-factor test.}
If the model predicts $\chi$ in Eq.~\eqref{eq:form_factor_template}, then
scattering data constrain the effective charge response.  The scale
$r_c\simeq\SI{1.409e-15}{\meter}$ cannot be interpreted naively as a
resolved electron radius; it is only observable through the predicted
coupling coefficient $\chi$.

\medskip
\noindent\textbf{(iii) Tension-ratio diagnostic.}
The ratio
\begin{equation}
  \frac{\Tc}{\Fgr}=\frac{4\alpha_g}{\tilde\alpha}
  \label{eq:tension_ratio_diagnostic}
\end{equation}
is a useful dimensionless diagnostic, but not by itself a falsifiable
prediction.  It becomes testable only after a concrete analogue-system
mapping identifies the measured mechanical tension and the analogue
causal-horizon scale.

% =============================================================
\section{Numerical Summary}
\label{sec:numerical}
% =============================================================

\begin{table}[h]
\centering
\caption{Numerical values of all closed quantities from CODATA 2022
inputs~\cite{MohrTaylor2022}.
The fourth column gives the propagated fractional uncertainty from $\Rinf$
(\num{1.9e-12}) where applicable; dashes indicate quantities with no
CODATA counterpart.}
\label{tab:numbers}
\begin{tabular}{llll}
\toprule
Quantity & Symbol & Value & Prop.\ unc.\ from $\Rinf$ \\
\midrule
Circ.\ speed   & $\vnorm$ & \SI{1.093\,845\,63e6}{\m\per\s}      & $\sim\num{1e-12}$ \\
Core radius    & $\rc$    & \SI{1.408\,970\,16e-15}{\meter}       & $\sim\num{1e-12}$ \\
$\tilde\alpha$ &          & \num{7.297\,352\,567e-3}              & $\sim\num{1e-12}$ \\
Circulation    & $\Gz$    & \SI{9.683\,619e-9}{\m\squared\per\s}  & — \\
Line tension   & $\Tc$    & \SI{29.053\,51}{\newton}               & — \\
Core density   & $\rhocore$& \SI{3.893\,4e18}{\kg\per\m\cubed}    & — \\
Geom.\ gate    & $4/\tilde\alpha$ & \num{548.14}               & $\sim\num{1e-12}$ \\
\bottomrule
\end{tabular}
\end{table}

% =============================================================
\section{Discussion}
\label{sec:discussion}
% =============================================================

The central result is a closure analysis, not a first-principles
derivation.  The Compton-core postulate fixes the relation between
rotation rate, speed, and radius.  The Rydberg anchor fixes the speed
only because it is equivalent to the fine-structure constant.  The line
tension and density then follow only after two further mechanical
closures are imposed.

This organization removes the circularity in the stronger claim that
$\alpha$ or $\rho_{\mathrm{core}}$ had been derived.  The paper instead
identifies exactly where new dynamics would be required.  The most
important missing piece is an independent derivation of
$\tilde\alpha=2\vnorm/c$ from a specified interaction functional.  The
second missing piece is an on-shell core-energy calculation yielding
$\mathcal{T}_c=m_ec^2/(2r_c)$.  The third is a stress-balance derivation
of the thin-wall Laplace law.

The background fluid density $\rhof$ remains open.  Its determination
requires an independent electromagnetic--circulatory coupling or another
observable normalization.  It is therefore deliberately excluded from
the claimed closure chain.

% =============================================================
\section{Conclusion}
\label{sec:conclusion}
% =============================================================

A Compton-core vortex-filament model becomes algebraically closed once
four ingredients are supplied: the Compton-core rotation postulate, the
spectroscopic speed closure, the core-diameter line-tension
normalization, and the thin-wall pressure-balance closure.  These give
\begin{align*}
  \vnorm  &= \sqrt{\pi\hbar cR_\infty/m_e},
  &\rc    &= \vnorm/\omega_C = r_e/2, \\
  \Gz     &= 2\pi\rc\vnorm,
  &\Tc    &= m_ec^2/(2\rc),\\
  \rhocore &= m_ec^2/(2\pi\vnorm^2\rc^3).
\end{align*}
The relations are dimensionally consistent and internally coherent, but
only the first is a structural postulate and the others depend on
explicit closures.  A future first-principles derivation must replace
those closures by a dynamical calculation of $\alpha$, the line tension,
the interfacial stress law, and at least one low-energy observable
coefficient.

% =============================================================
\begin{acknowledgments}
The author acknowledges helpful discussions on the distinction between
algebraic closure, empirical calibration, and first-principles
derivation.
\end{acknowledgments}

% =============================================================
\begin{thebibliography}{99}

\bibitem{Helmholtz1867}
H.~von Helmholtz,
On the integrals of the hydrodynamical equations which express
vortex-motion,
\textit{Philosophical Magazine} \textbf{33}, 485 (1867).

\bibitem{Kelvin1867}
Lord Kelvin (W.~Thomson),
On vortex atoms,
\textit{Proceedings of the Royal Society of Edinburgh} \textbf{6},
94 (1867).

\bibitem{MohrTaylor2022}
P.~J. Mohr, D.~B. Newell, B.~N. Taylor, and E.~Tiesinga,
CODATA recommended values of the fundamental physical constants: 2022,
arXiv:2409.03787 (2024).

\bibitem{Donnelly1991}
R.~J. Donnelly,
\textit{Quantized Vortices in Helium II},
Cambridge University Press, 1991.

\bibitem{Pethick2008}
C.~J. Pethick and H.~Smith,
\textit{Bose--Einstein Condensation in Dilute Gases}, 2nd ed.,
Cambridge University Press, 2008.
\doi{10.1017/CBO9780511802850}

\bibitem{FaddeevNiemi1997}
L.~D. Faddeev and A.~J. Niemi,
Stable knot-like structures in classical field theory,
\textit{Nature} \textbf{387}, 58 (1997).
\doi{10.1038/387058a0}

\bibitem{Barcelo2005}
C.~Barcel\'o, S.~Liberati, and M.~Visser,
Analogue gravity,
\textit{Living Reviews in Relativity} \textbf{8}, 12 (2005).
\doi{10.12942/lrr-2005-12}

\bibitem{Volovik2003}
G.~E. Volovik,
\textit{The Universe in a Helium Droplet},
Oxford University Press, 2003.

\bibitem{Saffman1992}
P.~G. Saffman,
\textit{Vortex Dynamics},
Cambridge University Press, 1992.
\doi{10.1017/CBO9780511624063}

\bibitem{Onsager1949}
L.~Onsager,
Statistical hydrodynamics,
\textit{Nuovo Cimento Supplement} \textbf{6}, 279 (1949).
\doi{10.1007/BF02780991}

\bibitem{Feynman1955}
R.~P. Feynman,
Application of quantum mechanics to liquid helium,
in \textit{Progress in Low Temperature Physics}, vol.~1,
ed.\ C.~J. Gorter, p.~17, North-Holland, 1955.

\bibitem{Gibbons2002}
G.~W. Gibbons,
The maximum tension principle in general relativity,
\textit{Foundations of Physics} \textbf{32}, 1891 (2002).
\doi{10.1023/A:1022370717626}

\bibitem{Schiller2006}
C.~Schiller,
Simple derivation of minimum length, minimum dipole moment, and
lack of space-time continuity,
\textit{International Journal of Theoretical Physics} \textbf{45},
221 (2006).
\doi{10.1007/s10773-005-4835-2}

\bibitem{Karshenboim2005}
S.~G. Karshenboim,
Precision physics of simple atoms: QED tests, nuclear structure, and
fundamental constants,
\textit{Physics Reports} \textbf{422}, 1 (2005).
\doi{10.1016/j.physrep.2005.08.008}

\bibitem{Bethe1947}
H.~A. Bethe,
The electromagnetic shift of energy levels,
\textit{Physical Review} \textbf{72}, 339 (1947).
\doi{10.1103/PhysRev.72.339}

\bibitem{Batchelor1967}
G.~K. Batchelor,
\textit{An Introduction to Fluid Dynamics},
Cambridge University Press, 1967.
\doi{10.1017/CBO9780511800955}

\bibitem{PDG2022}
R.~L. Workman et al.\ (Particle Data Group),
Review of particle physics,
\textit{Progress of Theoretical and Experimental Physics}
\textbf{2022}, 083C01 (2022).
\doi{10.1093/ptep/ptac097}

\end{thebibliography}

\end{document}
